\chapter{Introduction}

\section{Background}
The field of marine structural engineering has historically been characterized by conservative design principles, driven by the harsh and unpredictable nature of the ocean environment. Naval architects and marine engineers are tasked with designing vessels and offshore structures that must withstand complex, dynamic loads including wave slamming, vibration from propulsion systems, and extreme weather conditions. Consequently, the design of marine components---ranging from massive hull girders to auxiliary equipment mounting brackets---has traditionally prioritized robust structural integrity over mass optimization. However, in the modern maritime industry, there is an escalating demand for high-efficiency vessels. Stricter international emissions regulations, fluctuating fuel costs, and the economic imperative to maximize payload capacity have made weight reduction a critical objective in ship design. Every kilogram of unnecessary structural mass contributes to increased fuel consumption and reduced operational efficiency over a vessel's multi-decade service life.

Despite the widespread adoption of Computer-Aided Design (CAD) and Finite Element Analysis (FEA) software, the standard design workflow remains largely manual and iterative. An engineer typically develops an initial geometric concept, applies boundary conditions and loads in an FEA environment, and solves for stresses and deformations. If the design fails to meet the required safety margins, the geometry is fortified; if it passes by a wide margin, the engineer might attempt to remove material to save weight. This process is repeated until a satisfactory design is achieved. While CAD and FEA tools have dramatically improved the accuracy of engineering analysis, the human-in-the-loop iteration process severely restricts the design space that can be explored. 

Algorithmic optimization offers a paradigm shift. By coupling parametric CAD models with external numerical optimization routines, the iterative design loop can be automated. A computational pipeline can evaluate dozens, hundreds, or even thousands of design variants in the time it takes a human engineer to evaluate a handful. Gradient-free optimization algorithms, such as the Nelder-Mead simplex method, are particularly well-suited for this application because they do not require analytical derivatives of the FEA solver, treating the CAD/FEA environment as a black box. This project introduces the Nelder-Mead Marine Optimizer (NEMO), an automated, FEA-driven design optimization pipeline designed to bridge the gap between parametric CAD and external optimization logic, specifically tailored for marine structural components.

\section{Problem Statement}
The core problem addressed by this project is the profound inefficiency of manual, trial-and-error design iteration in structural engineering, particularly within the marine sector. The financial and temporal costs of manual iteration are substantial. Human engineers are slow at manually updating geometry, re-meshing models, running simulations, and interpreting results. As a consequence, design iterations are expensive in terms of labor hours, forcing engineering teams to severely limit the number of variants they can realistically evaluate before project deadlines. 

Due to these temporal constraints, engineers typically settle for designs that are merely ``good enough.'' Once a design satisfies the requisite structural safety requirements (e.g., a minimum factor of safety), the iteration process often ceases. This premature convergence leaves significant performance on the table. In the context of marine engineering, this translates directly to excess material usage, higher manufacturing costs, and unnecessary vessel deadweight. For service-critical applications, particularly moving parts or widely replicated fittings, this excess mass cascades into negative impacts on fuel consumption and lifecycle efficiency. 

Furthermore, human intuition is inherently limited when navigating multi-dimensional design spaces. An engineer might intuitively understand that thickening a baseplate reduces stress, but finding the exact, optimal combination of baseplate thickness, rib height, fillet radius, and overall length that minimizes mass while simultaneously satisfying both stress and deflection constraints is mathematically complex. The lack of automated optimization pipelines in standard, accessible engineering workflows prevents the discovery of genuinely optimal (minimum-mass, constraint-satisfying) structural designs within normal project timelines. 

\section{Objectives}

\subsection{Primary Objective}
The primary objective of this project is to develop, implement, and validate the Nelder-Mead Marine Optimizer (NEMO): an automated, FEA-driven software pipeline that couples Autodesk Fusion 360 (parametric CAD and FEA) with an external Python-based optimization routine to systematically search for and converge upon a minimum-mass design of a marine structural component while strictly satisfying predefined structural safety constraints.

\subsection{Specific Objectives}
To achieve the primary objective, the following specific objectives will be met:
\begin{enumerate}
    \item To define a robust parametric CAD model of a marine structural component, exposing key geometric variables (such as baseplate dimensions, rib profiles, and fillet radii) for programmatic manipulation.
    \item To establish a pre-validated Finite Element Analysis (FEA) study template, incorporating realistic marine load cases (including static weight and dynamic amplification factors) and boundary conditions.
    \item To engineer a reliable, file-based handshake architecture between the Fusion 360 embedded Python API and an external Python execution environment, enabling bidirectional communication of parameter requests and FEA results.
    \item To implement a penalty-constrained gradient-free optimization algorithm (Nelder-Mead) capable of navigating the design space to minimize mass while penalizing constraint violations (e.g., factor of safety and maximum deflection).
    \item To quantify the benefits of algorithmic optimization by benchmarking the automated pipeline's final design against a manually-iterated baseline component in terms of total mass reduction and structural integrity.
\end{enumerate}

\section{Significance}
The significance of this project lies in making parameterized design optimization more accessible within standard engineering workflows. The transparent, file-based pipeline combines the Fusion 360 in-app API with a tested pure-Python optimizer and traceable validation records.

From a marine engineering perspective, the benefits of such algorithmic optimization are substantial. Marine structures are characterized by the repetition of numerous standardized components, such as brackets, stiffeners, and mounting foundations. An optimized algorithm can systematically shave off marginal percentages of mass from these components. Multiplied across the hundreds of similar fittings on a single vessel, this automated discovery of minimum-mass geometries translates into tangible reductions in material costs during fabrication and meaningful improvements in fuel efficiency and payload capacity over the vessel's lifetime. Furthermore, the pipeline relieves engineers from the tedious, repetitive tasks of manual CAD/FEA updates, allowing human expertise to be reallocated toward high-level system design and complex problem-solving.

\section{Scope and Limitations}
The scope of this project encompasses the development of the software pipeline architecture, the integration of the CAD/FEA environment with the Python optimizer, and the demonstration of the system's efficacy. 

The proposal is limited to a marine equipment mounting bracket supporting a small pump or auxiliary unit. A lifting padeye and a stabilizer structure are possible future case studies, but their design, analysis, and validation are outside the proposed bracket campaign.

Several aspects are explicitly out of scope for this phase of the project. Physical fabrication and experimental structural testing of the optimized designs will not be conducted; validation relies on the mathematical convergence of the simulated FEA results. Furthermore, the project employs a simplified standard static-plus-dynamic-factor approximation for the load case, rather than deriving loads from full Computational Fluid Dynamics (CFD) or transient seakeeping simulations. Optimization is strictly single-objective (minimizing mass subject to safety constraints) and does not currently evaluate multi-objective trade-offs such as manufacturability cost or fatigue life. 

A primary limitation of the current implementation is its reliance on sequential, localized CAD regenerations and FEA solves, bounding the evaluation speed by the local machine's computational hardware and the software's API overhead. Additionally, while the file-based handshake ensures robustness and simplicity in bypassing the limitations of the CAD software's embedded interpreter, it introduces minor latency due to polling loops, which may marginally increase total optimization time compared to deeply integrated native routines.
